% Generated draft: requires mathematical and source review.
\documentclass[11pt]{article}
\usepackage[margin=25mm]{geometry}
\usepackage{fontspec}
\setmainfont{Latin Modern Roman}
\usepackage{amsmath,amssymb,mathtools,unicode-math}
\setmathfont{Latin Modern Math}
\usepackage{enumitem,needspace,xurl,hyperref}
\setlength{\parindent}{0pt}
\setlength{\parskip}{5pt}
\setlength{\emergencystretch}{3em}
\newcommand{\N}{\mathbb{N}}
\newcommand{\Z}{\mathbb{Z}}
\newcommand{\Q}{\mathbb{Q}}
\newcommand{\R}{\mathbb{R}}
\newcommand{\C}{\mathbb{C}}
\newcommand{\F}{\mathbb{F}}
\newcommand{\A}{\mathbb{A}}
\newcommand{\PP}{\mathbb{P}}
\newcommand{\E}{\mathbb{E}}
\newcommand{\eps}{\varepsilon}
\newcommand{\dd}{\,\mathrm{d}}
\begin{document}
\section*{UnsolvedMath ID: 20003156}
\textbf{Source problem code:} catalogue entry
\subsection{Definitions and mathematical statement}
Let a group $G$ act isometrically on a Riemannian manifold $M$, with fixed orbit map $\pi:M\to M/G$. The source asks whether a changed metric on the orbit space is the quotient of some invariant metric on the same $M$. It does not specify the acting group, quotient regularity or whether any curvature bound must survive. The aid uses proper compact-Lie-group actions and smooth invariant Riemannian lifts, without asserting curvature preservation.

For a constant-orbit-type action, $M/G$ is a smooth quotient in the relevant setting. Choose a $G$-invariant horizontal complement to the vertical tangent bundle and a positive vertical metric. A prescribed smooth base metric $h$ lifts by taking $\pi^*h$ horizontally, the chosen metric vertically and the two bundles orthogonal. The aid also attributes to Mendes an affirmative lifting theorem for smooth orbifold metrics on orbifold quotients, including polar actions. Precise orbifold hypotheses and the full citation are absent.

An arbitrary ordinary boundary-smooth metric on a singular quotient need not lift. For the rotation action $\operatorname{SO}(2)\curvearrowright\R^2$, fix the quotient coordinate $r=\sqrt{x^2+y^2}$ and let the desired metric on $[0,\infty)$ be $q(r)\,\dd r^2$, $q>0$. The aid's exact criterion is
\[
\text{a smooth invariant exact lift exists}
\quad\Longleftrightarrow\quad q(r)=Q(r^2)
\]
for a smooth positive $Q$ on the squared-radius half-line, equivalently $q$ has a smooth even extension at $0$. Sufficiency is realized by $Q(x^2+y^2)(\dd x^2+\dd y^2)$. Thus $(1+r)\,\dd r^2$ has no such lift through the fixed orbit map. It is not a smooth orbifold metric at the reflection point in the needed sense, so this does not contradict the orbifold theorem. Changing the quotient coordinate by an arbitrary isometry is not part of the exact-lift question.

The aid's local-to-global method glues inverse metrics. Suppose a saturated open cover $\{U_i\}$ of $M$ has smooth $G$-invariant local metrics $g_i$ inducing the same prescribed quotient metric under the fixed local orbit maps. With a smooth invariant basic partition of unity $\{\rho_i\}$ subordinate to that cover, set
\[
g^{-1}=\sum_i\rho_i g_i^{-1},\qquad
g=\left(\sum_i\rho_i g_i^{-1}\right)^{-1}.
\]
The sum is a positive cometric. On covectors pulled back from each quotient stratum, every local cometric is the same prescribed quotient cometric, so the basic sum retains it. Under the stated exact quotient and stratified-length compatibility assumptions, the aid asserts that $g$ induces the prescribed quotient metric globally. This supplies a sufficient reduction to exact local lifts in slice models. It does not say that every singular quotient length metric has such local lifts or that averaging preserves an additional sectional-curvature condition.
\subsection{Short English statement}
Metric lifting depends on the quotient smooth structure. The supplied aid gives regular and orbifold positive cases, an even-jet obstruction at a rotational fixed point, and a local-to-global construction that glues exact lifts through inverse metrics.
\subsection{Sources}
\begin{itemize}
\item Local UnsolvedMath registry, permanent ID $20003156$, source code \texttt{catalogue entry}; canonical \texttt{aim-topology-notes.json}, \texttt{notes[243]}.
\item AIM workshop \emph{Manifolds with nonnegative sectional curvature}, Problem 31: \url{https://aimath.org/WWN/nnsectcurvature/nnsectcurvature.pdf}. The supplied verification reports agreement of the database text with the source wording. Original remarks and literature field are empty.
\item Mendes is named in the aid for lifting smooth orbifold metrics, without an authoritative full bibliographic entry or exact hypotheses.
\item Research attempt 1: \texttt{partial\_result}, status \texttt{partially\_solved}, classification \texttt{PARTIAL-PROGRESS}, rigor \texttt{proved}, novelty confidence low. Its even-jet and cometric-gluing criteria are machine-generated aid, not peer reviewed, and require independent verification before citation.
\end{itemize}
\end{document}
